\section{Four Dogmas, One Economy}

The Catholic Church counts four truths about the Blessed Virgin Mary among
its dogmas: that she is truly Mother of God (\term{Theotokos}); that she is
ever-virgin---before, in, and after the birth of her Son; that she was
conceived immaculate, preserved from all stain of original sin from the
first instant of her conception; and that, her earthly life completed, she
was assumed body and soul into heavenly glory. This reference states, for
each of the four, exactly what was defined, by whom, with what authority,
and from what scriptural and traditional roots; what the definition does
not assert; and how the Church presently teaches it.

A dogma, in the Church's strict usage, is a truth contained in the word of
God, written or handed down, which the Church proposes as divinely
revealed---whether by a solemn judgment or by her ordinary and universal
magisterium---and which therefore calls for the assent of divine and
Catholic faith. The First Vatican Council fixed that criterion in its
constitution \work{Dei Filius} (1870).\endnote{\work{Dei Filius}, ch.~3
(DS 3011), checked in the English presentation of the First Vatican
Council's decrees at
\url{https://www.papalencyclicals.net/councils/ecum20.htm}, 2026-07-25.
All web loci cited in these notes were read on 2026-07-25 unless another
date is stated; full bibliographic detail and witness identifications are
in the References. Denzinger numbers are cited as DS in the revised
numbering; where a claim depends on an online Denzinger witness, the
References identify it.} Not every true or authoritative statement about
Mary is a dogma: the four treated here are the ones commonly so identified
because each meets that strict standard. Their modes of proposal differ,
and the differences matter.

\dogmathesis{The four Marian dogmas confess one work of God, not four
privileges of Mary considered alone. Because the eternal Son truly became
man from her, she is Mother of God; because her whole self was given to
that motherhood, she is ever-virgin; because the Redeemer's merits reached
her first and most perfectly, she was preserved from original sin; because
her body shared the Redeemer's victory, she was assumed into glory. Each
dogma was defined in a different mode---two in the ancient conciliar and
traditional inheritance, two by solemn papal definition---and each leaves
identified questions deliberately open. Everything defined is received:
nothing in the four makes Mary a source of grace independent of Christ,
the sole Redeemer and Mediator.}

The two ancient dogmas were fixed in the Christological crisis of the
fifth to seventh centuries: the divine motherhood at the Council of
Ephesus (431) and in the councils that received it, the perpetual
virginity through the Church's universal tradition, confessed in creeds
and councils and given an exact threefold formulation by the Lateran
Council of 649. The two modern dogmas were defined \term{ex cathedra} by
Pius IX in \work{Ineffabilis Deus} (1854) and Pius XII in
\work{Munificentissimus Deus} (1950)---the two acts universally cited as
exercises of the papal teaching authority that the First Vatican Council
would define in 1870. Denial of a dogma is not mere disagreement with a
theological school; both modern definitions attach a formal sanction to
willful denial, quoted below with their texts.

Method, throughout: the exact defined text is quoted from an identified
witness and set off typographically; the doctrinal history is traced to
the definition, not read backwards from it; scriptural and patristic
evidence is weighted honestly, with translation cares stated; and the
boundary between dogma, authoritative teaching, and theological opinion is
kept explicit, above all for the disputed cooperation titles treated in
the final chapter. Work-wide bounds, terminology, and qualifications are
gathered in the terminal appendix.
